% Appendix, from exam/README.md: the form of the written exam, its grade limits, what it covers and
% how it is marked; and from info/README.md, its sections "Examination" (how the labs and the exam
% are graded, and the final grade) and "Lärandemål" (where each learning outcome is taught and
% examined), with the lectures L01 to L12 given as the chapters they became.
%
% Deliberate differences: everything that belongs to one run of the course is left out, the school
% and its routines, the dates, and the lecture the exam is written in ("i L12" is "i slutet av
% kursen"). info/README.md's own paragraph on the written exam repeats exam/README.md's, and is
% folded into it. The exam itself and its solutions are the teacher's and are not in the book.

\chapter{Det skriftliga provet}
\label{app:exam}

Kursen avslutas med ett skriftligt prov, efter kapitel~\ref{c12:ch}. Här står hur provet är upplagt,
vad det omfattar och hur det bedöms, och hur det tillsammans med laborationerna ger kursens betyg.

Ett övningsprov finns att träna på: frågorna står i bilaga~\ref{app:practice}, och
lösningsförslagen, med poäng per deluppgift, i bilaga~\ref{app:practice:answers}. Gör det under två
timmar, innan du tittar i lösningarna. Själva provet och dess lösningar är lärarens och finns inte
i boken.

\section*{Provets form}
\markright{Provets form}
\begin{itemize}
  \item \textbf{Två timmar}, individuellt, i slutet av kursen.
  \item \textbf{Hjälpmedel:} penna, suddgummi och linjal. Ingen miniräknare. Ett referensblad med
    AVR-instruktionerna, bilaga~\ref{app:avr}, delas ut tillsammans med provet; det är samma blad
    som använts under kursen.
  \item \textbf{40 poäng}, fördelade på sex frågor. Varje fråga har deluppgifter, och poängen står
    vid varje deluppgift.
  \item Svaren skrivs på separata papper, med en ny fråga på ett nytt papper.
\end{itemize}

\subsection*{Betygsgränser}

\begin{narrowtable}{@{}ll@{}}
  \thead{Betyg} & \thead{Poäng} \\
  \midrule
  U & under 20 \\
  G & 20--29 \\
  VG & 30--40 \\
\end{narrowtable}

\noindent Slutbetyget på kursen bygger på provet och laborationerna tillsammans; se \enquote{Provet
och laborationerna} nedan.

\section*{Vad provet omfattar}
\markright{Vad provet omfattar}
Provet omfattar hela kursen utom sekvensnät och vippor, som förekommer bara som en del av
mikrodatorns uppbyggnad. Frågorna är desamma till form och poäng på varje prov:

\begin{booktable}{@{}p{2.6em}Lp{3.6em}Lr@{}}
  \thead{Fråga} & \thead{Område} & \thead{Kapitel} & \thead{Lärandemål} & \thead{Poäng} \\
  \midrule
  1 & Talsystem och koder & 1 & beskriva talsystem; omvandla talvärden & 8 \\
  2 & Grindar och boolesk algebra & 2 & förklara grindar; beskriva logisk algebra och förenkling &
    8 \\
  3 & Karnaughdiagram & 4 & använda Karnaughdiagram för minimering & 6 \\
  4 & Analys av ett kombinatoriskt nät, grindnät eller kontaktnät & 2--5 & analysera digitala
    kombinatoriska nät & 6 \\
  5 & Mikrodatorns uppbyggnad & 6--7 & redogöra för mikrodatorns uppbyggnad & 4 \\
  6 & Assemblerprogrammering & 7--10 & använda assemblerspråk för att programmera en mikrodator &
    8 \\
  \midrule
    & & & & \textbf{40} \\
\end{booktable}

\noindent Lärandemålen om att syntetisera nät med IC-kretsar och kontakter, och att styra
styrobjekt, examineras i laborationerna, där de hör hemma.

Assemblerfrågorna gäller ATmega328P med klockfrekvensen 16~MHz, och de instruktioner som står på
referensbladet. Ingen fråga kräver att en instruktionskod eller en cykeltid kan utantill; det som
behövs står på bladet.

\section*{Hur provet bedöms}
\markright{Hur provet bedöms}
\begin{itemize}
  \item \textbf{Metoden ger poäng.} En korrekt uträkning med ett slarvfel är värd mer än ett rätt
    svar utan uträkning. Ett svar utan uträkning kan få full poäng bara där frågan inte ber om
    någon.
  \item \textbf{Följdfel kostar en gång.} Om en deluppgift bygger på svaret i en tidigare, och det
    svaret var fel, bedöms den senare deluppgiften utifrån det felaktiga svaret. Ett fel i en
    sanningstabell som förs vidare till ett Karnaughdiagram kostar poäng i tabellen, inte i
    diagrammet.
  \item \textbf{Talsystemet ska synas.} Ett svar som \code{2C} utan skrivsätt är tvetydigt och ger
    inte full poäng. Skriv \code{0x2C}, \code{0010 1100} eller 44.
  \item \textbf{Program bedöms efter vad de gör, inte efter stavningen.} Ett program som skulle
    fungera men har ett felstavat registernamn eller ett saknat kommatecken förlorar ingenting. Ett
    program som läser fel bit, eller glömmer att en knapp är aktivt låg, förlorar de poäng den delen
    gäller.
  \item \textbf{Namngivna fällor.} Några deluppgifter finns till för att se om ett vanligt misstag
    undviks, till exempel en flagga som en instruktion inte ändrar, eller en Karnaughaxel i binär
    ordning i stället för Graykod. Lösningsförslagen säger var fällorna finns, och ett svar som går
    i en fälla förlorar de poäng som fällan gäller, inte hela deluppgiftens. Lösningsförslagen till
    övningsprovet, i bilaga~\ref{app:practice:answers}, visar hur det ser ut.
\end{itemize}

\section*{Provet och laborationerna}
\markright{Provet och laborationerna}
\begin{itemize}
  \item \textbf{Labb~1} och \textbf{Labb~2} bedöms med \textbf{U/G}.
  \item \textbf{Labb~3} bedöms med \textbf{U/G/VG}.
  \item Det \textbf{skriftliga provet}, två timmar i slutet av kursen, bedöms med \textbf{U/G/VG}.
    \textbf{G} kräver 20 poäng av 40, och \textbf{VG} kräver 30.
\end{itemize}

\subsection*{Laborationerna}
En laboration är godkänd när samtliga uppgifter för godkänt är redovisade för läraren, och båda i
paret kan förklara lösningen. Redovisning sker under laborationspassen. Handledningarna finns i
bilaga~\ref{app:lab1}, \ref{app:lab2} och~\ref{app:lab3}.

\begin{booktable}{@{}p{5.6em}LL@{}}
  \thead{Laboration} & \thead{Krav för G} & \thead{Krav för VG} \\
  \midrule
  Labb~1 & Uppgift 1--8 & -- \\
  Labb~2 & Labb~2-1 uppgift 1--4 och Labb~2-2 uppgift 1--3 & -- \\
  Labb~3 & Del 1.1--12, 14--16, 21 och slutuppgiften del A & Samtliga delar 1.1--23 och
    slutuppgiften del B \\
\end{booktable}

\subsection*{Slutbetyg}

\begin{booktable}{@{}p{5.6em}L@{}}
  \thead{Betyg} & \thead{Krav} \\
  \midrule
  G & Samtliga laborationer godkända och godkänt skriftligt prov. \\
  VG & Samtliga laborationer godkända, VG på Labb~3 och VG på det skriftliga provet. \\
\end{booktable}

\noindent För godkänt betyg krävs att alla kursens lärandemål är uppnådda. Omprov och
restlaborationer genomförs enligt skolans rutiner.

\needspace{20\baselineskip}
\section*{Lärandemålen}
\markright{Lärandemålen}
Var varje lärandemål i kursplanen undervisas, och var det examineras.

\begin{booktable}{@{}Lp{5.4em}p{9em}@{}}
  \thead{Lärandemål} & \thead{Kapitel} & \thead{Examineras} \\
  \midrule
  Förklara grindar och logiska funktioner & 2, 4 & Labb~1, Labb~2, prov \\
  Beskriva binära, hexadecimala och decimala talsystem & 1 & Prov \\
  Beskriva logisk algebra och förenklingar av booleska uttryck & 2, 4 & Prov \\
  Redogöra för mikrodatorns uppbyggnad på blockschemanivå & 6, 7 & Prov \\
  Omvandla talvärden mellan olika talsystem & 1, 7, 8 & Prov, Labb~3 \\
  Syntetisera digitala kombinatoriska nät med IC-kretsar och kontakter & 2--5 & Labb~1, Labb~2 \\
  Använda Karnaughdiagram för analys och minimering av booleska uttryck & 4 & Labb~2, prov \\
  Använda assemblerspråk för att programmera en mikrodator & 7--10 & Labb~3, prov \\
  Programmera mikrodatorn till att kontrollera styrobjekt & 8, 10, 11 & Labb~3 \\
  Analysera digitala kombinatoriska nät & 2--5 & Labb~1, Labb~2, prov \\
\end{booktable}

\noindent Sekvensnät och vippor (kapitel~\ref{c6:ch}) ingår inte i kursplanens lärandemål. De tas
med för att mikrodatorn annars blir en svart låda: register, programräknare och minne är vippor,
och det är den kopplingen som gör blockschemat begripligt. På provet förekommer vippor bara som en
del av mikrodatorns uppbyggnad.
